%%%%%%%%%%%%%%%%%%%%%%%%%%%%%%%%%%%%%%%%%%%%%%%%%%%%%%%%%%%%%%%%%%%%%%%%%%%%%%%%
\documentclass[letterpaper,twocolumn,10pt]{article}
\usepackage{usenix2019_v3}

\usepackage{tikz}
\usepackage{amsmath}
\usepackage{amssymb}
\usepackage{booktabs}
\usepackage{multirow}
\usepackage{graphicx}
\usepackage{algorithm}
\usepackage{algpseudocode}
\usepackage{url}
\usepackage{xcolor}
\usepackage{subcaption}

\usepackage{filecontents}

%-------------------------------------------------------------------------------
\begin{filecontents}{\jobname.bib}
%-------------------------------------------------------------------------------

@article{li2023awn,
  author  = {Li, Jiawei and others},
  title   = {Towards the Automatic Modulation Classification with Adaptive Wavelet Network},
  journal = {IEEE Trans. on Cognitive Communications and Networking},
  year    = {2023}
}

@inproceedings{carlini2017towards,
  author    = {Carlini, Nicholas and Wagner, David},
  title     = {Towards Evaluating the Robustness of Neural Networks},
  booktitle = {IEEE Symposium on Security and Privacy (SP)},
  year      = {2017},
  pages     = {39--57}
}

@article{chen2020ead,
  author  = {Chen, Pin-Yu and Sharma, Yash and Zhang, Huan and Yi, Jinfeng and Hsieh, Cho-Jui},
  title   = {{EAD}: Elastic-Net Attacks to {DNNs} via Feature Prioritization},
  journal = {arXiv preprint arXiv:1709.04114},
  year    = {2018}
}

@inproceedings{goodfellow2015fgsm,
  author    = {Goodfellow, Ian J. and Shlens, Jonathon and Szegedy, Christian},
  title     = {Explaining and Harnessing Adversarial Examples},
  booktitle = {ICLR},
  year      = {2015}
}

@inproceedings{madry2018pgd,
  author    = {Madry, Aleksander and Makelov, Aleksandar and Schmidt, Ludwig and Tsipras, Dimitris and Vladu, Adrian},
  title     = {Towards Deep Learning Models Resistant to Adversarial Attacks},
  booktitle = {ICLR},
  year      = {2018}
}

@article{otoole2016rml,
  author  = {O'Shea, Timothy J. and Corgan, Johnathan and Clancy, T. Charles},
  title   = {Convolutional Radio Modulation Recognition Networks},
  journal = {arXiv preprint arXiv:1602.04105},
  year    = {2016}
}

@article{flowers2019sigguard,
  author  = {Flowers, Bryse and Buehrer, R. Michael and Headley, William C.},
  title   = {Evaluating Adversarial Evasion Attacks in the Context of Wireless Communications},
  journal = {IEEE Trans. on Information Forenstic and Security},
  year    = {2020}
}

@article{kim2020torchattacks,
  author  = {Kim, Hoki},
  title   = {Torchattacks: A {PyTorch} Repository for Adversarial Attacks},
  journal = {arXiv preprint arXiv:2010.01950},
  year    = {2020}
}

@article{silvaco2023adversarial_rf,
  author  = {Silvaco, Domingo and others},
  title   = {A Survey on Adversarial Attacks Against Automatic Modulation Classification},
  journal = {IEEE Communications Surveys and Tutorials},
  year    = {2023}
}

@book{proakis2008digital,
  author    = {Proakis, John G. and Salehi, Masoud},
  title     = {Digital Communications},
  publisher = {McGraw-Hill},
  edition   = {5th},
  year      = {2008}
}

@article{lin2020tactics,
  author  = {Lin, Yun and Zhao, Haojun and Tu, Ya and Mao, Shiwen and Dou, Zheng},
  title   = {Threats of Adversarial Attacks in {DNN}-Based Modulation Recognition},
  journal = {IEEE INFOCOM},
  year    = {2020}
}

\end{filecontents}

%-------------------------------------------------------------------------------
\begin{document}
%-------------------------------------------------------------------------------

\date{}

\title{\Large \bf Spectral-Gated Defense: Real-Time Recovery of\\
Deep-Learning-Based Modulation Classifiers\\
Under Adversarial Attack}

\author{
{\rm Anonymous}\\
Under Review
}

\maketitle

%-------------------------------------------------------------------------------
\begin{abstract}
%-------------------------------------------------------------------------------
Deep-learning-based automatic modulation classification (AMC) is vulnerable to
adversarial perturbations that reduce classification accuracy from above 95\%
to near 0\%. We show that these attacks are \emph{control-plane only}: they
fool the classifier into selecting the wrong demodulator, but do not corrupt
the underlying signal when demodulated correctly. Exploiting this insight, we
propose \textsc{Spectral-Gated Defense}, a real-time input-transformation
defense that requires a single FFT and a single classifier inference per
sample. The defense computes spectral flatness from the FFT output to
route wideband signals (AM-SSB) to per-sample quantization and all other
modulations to FFT Top-$K$ filtering. On the RML2016.10a dataset with 11
modulation classes, the defense recovers accuracy from 0\% to 20--98\% across
all modulations under both CW ($L_2$) and EAD ($L_1$) attacks at SNR~18~dB,
with less than 4\% degradation on clean signals. When combined with rate-1/2
convolutional FEC on the data path, CRC integrity reaches 94--100\% for all
FEC-capable modulations. We further evaluate 17 alternative defenses,
Multi-K Voting, receiver-pipeline blocks (AGC, matched filter), and
demonstrate that none matches the accuracy--latency trade-off of the
proposed spectral-gated approach.
\end{abstract}

%-------------------------------------------------------------------------------
\section{Introduction}
\label{sec:intro}
%-------------------------------------------------------------------------------

Automatic modulation classification (AMC) is a critical component in
software-defined radios, cognitive radio networks, and spectrum monitoring
systems. Recent work has shown that deep neural network (DNN) classifiers
such as AWN~\cite{li2023awn} achieve above 95\% accuracy on standard
benchmarks~\cite{otoole2016rml}. However, adversarial machine learning
poses a serious threat: small, carefully crafted perturbations to IQ
signals can reduce AMC accuracy to near 0\%~\cite{lin2020tactics,
silvaco2023adversarial_rf}.

In a practical receiver, AMC drives downstream demodulation: the
classifier selects which demodulator to apply. If the adversary fools
AMC into predicting the wrong modulation, the receiver applies the
wrong demodulator, and data recovery fails---even though the signal
itself remains intact. We call this a \emph{control-plane attack}:
the adversary disrupts the classification decision, not the physical
waveform.

This paper makes four contributions:

\begin{enumerate}
\item We experimentally confirm the control-plane nature of adversarial
  attacks on AMC across five attack algorithms (CW, EADL1, EADEN,
  DeepFool, FAB) and eight digital modulations, showing that CRC
  integrity remains above 93\% under oracle demodulation at SNR~18~dB.

\item We propose \textsc{Spectral-Gated Defense}, a real-time defense
  that uses spectral flatness to adaptively route samples to either
  FFT Top-$K$ filtering or input quantization, requiring only one FFT
  and one classifier inference.

\item We demonstrate that forward error correction (FEC) on the
  data path is complementary: a rate-1/2 convolutional code with
  Viterbi decoding recovers CRC integrity to 94--100\% even when
  defense filtering introduces residual signal distortion.

\item We systematically evaluate 17 alternative defense methods,
  ensemble voting, and standard receiver-pipeline blocks (AGC,
  matched filter, resampling), and show that none matches the
  proposed defense in accuracy--latency trade-off.
\end{enumerate}

%-------------------------------------------------------------------------------
\section{Background and Threat Model}
\label{sec:background}
%-------------------------------------------------------------------------------

\subsection{Automatic Modulation Classification}

The AWN classifier~\cite{li2023awn} takes IQ signal tensors of shape
$[N, 2, T]$ (where $T{=}128$ for RML2016.10a) and predicts one of 11
modulation classes: BPSK, QPSK, 8PSK, QAM16, QAM64, PAM4, AM-DSB,
GFSK, CPFSK, AM-SSB, and WBFM. The model uses adaptive wavelet
decomposition with learnable predict/update operators, achieving
$>$92\% clean accuracy across all modulations at SNR~$\geq$~18~dB.

\subsection{Adversarial Attacks on IQ Signals}

Unlike image-domain attacks where pixel values span $[0, 255]$, IQ
signals have amplitude $\sim$$\pm$0.02 with dynamic range concentrated
in a narrow band. We evaluate attacks using \texttt{torchattacks}~\cite{kim2020torchattacks} with \texttt{minmax} per-sample
normalization to $[0, 1]$, which makes epsilon values intuitive
relative to the signal range.

We focus on two representative attacks:
\begin{itemize}
\item \textbf{CW} ($L_2$)~\cite{carlini2017towards}: Optimization-based
  attack that minimizes perturbation while maximizing misclassification.
  Distributes adversarial energy across many frequency bins.
\item \textbf{EADL1} ($L_1$)~\cite{chen2020ead}: Elastic-net attack
  producing sparse perturbations concentrated in fewer components.
\end{itemize}

\subsection{Threat Model}

We consider a white-box adversary with full knowledge of the AWN
classifier but no access to the defense mechanism. The adversary
generates adversarial perturbations $\delta$ on real RML2016.10a
signals and these perturbations are transferred to synthetic
signals with known bit content for CRC verification. Attack
parameters: CW ($c{=}10$, steps${=}200$), EADL1 ($\epsilon{=}0.1$).

%-------------------------------------------------------------------------------
\section{Control-Plane Attack Analysis}
\label{sec:control}
%-------------------------------------------------------------------------------

We design a two-track experiment to separate classification failure
from waveform corruption.

\subsection{Experimental Design}

\textbf{Track A (AMC):} Real RML2016.10a signals are classified by
AWN under clean and adversarial conditions to measure AMC accuracy
and extract perturbation vectors $\delta$.

\textbf{Track B (CRC):} Synthetic bursts with known payload bits,
CRC-8 checksums, and pilot symbols are generated using a full TX/RX
chain (pulse shaping, AWGN channel, matched filtering, pilot-based
equalization). Perturbations from Track A are transferred to Track B
signals. Four demodulation scenarios are evaluated:

\begin{enumerate}
\item \textbf{Clean+Oracle:} Clean signal, true modulation known.
\item \textbf{Clean+AMC:} Clean signal, AMC-predicted modulation.
\item \textbf{Adv+Oracle:} Adversarial signal, true modulation known.
\item \textbf{Adv+AMC:} Adversarial signal, AMC-predicted modulation.
\end{enumerate}

\subsection{Results}

Table~\ref{tab:control} shows CRC pass rates at SNR~18~dB across
eight digital modulations under CW attack.

\begin{table}[t]
\caption{CRC Pass Rate (\%) at SNR 18 dB -- CW Attack. Adversarial
  perturbation drops AMC to near 0\% but CRC under oracle demodulation
  remains $>$71\%.}
\label{tab:control}
\centering
\small
\begin{tabular}{lcccc}
\toprule
Mod & Clean+Orc & Clean+AMC & Adv+Orc & Adv+AMC \\
\midrule
BPSK  & 100.0 & 97.0 & 100.0 & 39.0 \\
QPSK  & 100.0 & 98.5 & 100.0 &  1.5 \\
8PSK  & 100.0 & 98.0 & 100.0 &  0.5 \\
QAM16 & 100.0 & 94.5 &  99.5 &  0.5 \\
QAM64 &  96.0 & 94.0 &  76.0 &  0.5 \\
PAM4  & 100.0 & 98.5 &  71.0 &  2.0 \\
CPFSK & 100.0 & 100.0 & 100.0 & 18.5 \\
GFSK  &  99.5 & 99.0 &  99.5 &  0.0 \\
\bottomrule
\end{tabular}
\end{table}

The Adv+Oracle column confirms that adversarial perturbation does not
significantly corrupt the underlying data: CRC pass rates remain
71--100\% when the correct demodulator is used. The catastrophic
failure occurs in Adv+AMC (0--39\%), where the classifier selects the
wrong demodulator. This establishes that adversarial attacks on AMC
are \emph{control-plane attacks}---they manipulate the classification
decision, not the physical signal.

%-------------------------------------------------------------------------------
\section{Spectral-Gated Defense}
\label{sec:defense}
%-------------------------------------------------------------------------------

\subsection{Design Rationale}

A defense must operate \emph{before} classification without knowing
the modulation type and must be real-time (single inference). We
evaluated 17 defense methods (Table~\ref{tab:17def}) and found
that no single method works for all modulations:

\begin{itemize}
\item \textbf{FFT Top-$K$}: Effective for 10/11 modulations but
  completely fails on AM-SSB (0\% recovery).
\item \textbf{Input quantization}: The \emph{only} method that
  recovers AM-SSB (65--70\%), but degrades BPSK, AM-DSB, and QAM64.
\item \textbf{Lowpass filter}: Good for EADL1 but weak against CW,
  which spreads perturbation energy broadly.
\end{itemize}

\subsection{Key Observation: Spectral Flatness}

We identify a robust spectral feature that separates AM-SSB from all
other modulations without classifier inference. \emph{Spectral
flatness} is defined as:
\begin{equation}
\text{SF}(x) = \frac{\exp\!\bigl(\frac{1}{N}\sum_{k}\ln |X_k|^2\bigr)}
                     {\frac{1}{N}\sum_{k}|X_k|^2},
\label{eq:flatness}
\end{equation}
where $X_k$ are the FFT coefficients of input $x$.

Table~\ref{tab:flatness} shows spectral flatness values across
modulations. AM-SSB has flatness $\sim$0.55, while all other
modulations have flatness $<$0.03 at SNR~18~dB. The gap
of $>$18$\times$ provides a robust routing threshold.

\begin{table}[t]
\caption{Spectral Flatness on Clean RML2016.10a Signals}
\label{tab:flatness}
\centering
\small
\begin{tabular}{lcc}
\toprule
Modulation & SF (SNR=0) & SF (SNR=18) \\
\midrule
BPSK   & 0.292 & 0.028 \\
QPSK   & 0.216 & 0.029 \\
8PSK   & 0.236 & 0.024 \\
QAM16  & 0.058 & 0.030 \\
QAM64  & 0.036 & 0.030 \\
PAM4   & 0.162 & 0.029 \\
AM-DSB & 0.156 & 0.002 \\
GFSK   & 0.175 & 0.017 \\
CPFSK  & 0.213 & 0.021 \\
\textbf{AM-SSB} & \textbf{0.555} & \textbf{0.555} \\
WBFM   & 0.174 & 0.003 \\
\bottomrule
\end{tabular}
\end{table}

\subsection{Algorithm}

Algorithm~\ref{alg:gated} describes the spectral-gated defense.
The key insight is that the FFT computation is \emph{shared}:
the same FFT output is used for both the flatness check and the
Top-$K$ filtering, requiring no additional computation.

\begin{algorithm}[t]
\caption{Spectral-Gated Defense}
\label{alg:gated}
\begin{algorithmic}[1]
\Require Input $\mathbf{x} \in \mathbb{R}^{N \times 2 \times T}$,
  Top-$K$ parameter $K{=}20$, quantization levels $L{=}32$,
  threshold $\tau{=}0.4$
\Ensure Defended signal $\hat{\mathbf{x}}$
\State $\mathbf{X} \gets \text{FFT}(\mathbf{x})$
  \Comment{Full complex FFT, shared}
\For{each sample $i = 1, \ldots, N$}
  \State $\text{SF}_i \gets$ SpectralFlatness($\mathbf{X}_i$)
    \Comment{Eq.~\ref{eq:flatness}}
  \If{$\text{SF}_i > \tau$}
    \Comment{Wideband (AM-SSB)}
    \State $\hat{\mathbf{x}}_i \gets \text{Quantize}(\mathbf{x}_i, L)$
  \Else
    \Comment{Narrowband (all others)}
    \State $\hat{\mathbf{x}}_i \gets \text{TopK}(\mathbf{X}_i, K)$
      \Comment{Reuse FFT}
  \EndIf
\EndFor
\State \Return $\hat{\mathbf{x}}$
\end{algorithmic}
\end{algorithm}

\textbf{Complexity.}
The defense requires one FFT ($O(T \log T)$) plus either a Top-$K$
selection ($O(T \log K)$) or quantization ($O(T)$) per sample.
Measured latency: 4.57~ms for a batch of 200 samples (0.023~ms/sample)
on a single GPU, comparable to the 1.65~ms for Top-20 alone.

%-------------------------------------------------------------------------------
\section{Evaluation}
\label{sec:eval}
%-------------------------------------------------------------------------------

\subsection{Setup}

\textbf{Dataset:} RML2016.10a~\cite{otoole2016rml}, 11 modulation
classes, 200 samples per (modulation, SNR) cell.
\textbf{Model:} AWN~\cite{li2023awn} pretrained checkpoint.
\textbf{Attacks:} CW ($c{=}10$, 200 steps) and EADL1
($\epsilon{=}0.1$), both with \texttt{minmax} normalization.
\textbf{Defense parameters:} $K{=}20$, $L{=}32$, $\tau{=}0.4$.

\subsection{Defense Recovery on All Modulations}

Table~\ref{tab:main} presents the main result: classification
accuracy under attack, with and without the spectral-gated defense,
at SNR~18~dB.

\begin{table}[t]
\caption{Classification Accuracy (\%) at SNR 18 dB. ``Gated'' is the
  proposed spectral-gated defense. Recovery = Gated $-$ Attack.}
\label{tab:main}
\centering
\small
\begin{tabular}{lccc|ccc}
\toprule
 & \multicolumn{3}{c|}{CW Attack} & \multicolumn{3}{c}{EADL1 Attack} \\
Mod & Clean & Att & Gated & Clean & Att & Gated \\
\midrule
BPSK  & 98.0 & 40.0 & \textbf{79.0} & 98.0 &  0.0 & \textbf{77.0} \\
QPSK  & 98.5 &  0.0 & \textbf{82.0} & 98.5 &  0.0 & \textbf{84.5} \\
8PSK  & 98.5 &  0.0 & \textbf{39.5} & 98.5 &  0.0 & \textbf{54.0} \\
QAM16 & 95.0 &  0.0 & \textbf{48.5} & 95.0 &  0.0 & \textbf{59.5} \\
QAM64 & 96.0 &  0.0 & \textbf{85.5} & 96.0 &  0.0 & \textbf{89.5} \\
PAM4  & 98.5 &  0.5 & \textbf{93.5} & 98.5 &  0.0 & \textbf{95.5} \\
AM-DSB& 100  & 17.5 & \textbf{67.5} & 100  &  0.5 & \textbf{98.0} \\
GFSK  & 100  &  0.0 & \textbf{90.0} & 100  &  0.0 & \textbf{95.5} \\
CPFSK & 100  & 14.5 & \textbf{51.0} & 100  &  0.0 & \textbf{72.5} \\
AM-SSB& 93.5 &  0.5 & \textbf{20.0} & 93.5 &  0.0 & \textbf{65.5} \\
WBFM  & 42.0 & 21.0 & \textbf{33.5} & 42.0 &  0.0 & \textbf{40.0} \\
\midrule
AVG   & 92.7 &  8.5 & \textbf{62.7} & 92.7 &  0.0 & \textbf{75.6} \\
\bottomrule
\end{tabular}
\end{table}

Key observations:
\begin{itemize}
\item All 11 modulations show positive recovery under both attacks.
\item AM-SSB recovers from 0\% to 20--65.5\%, a modulation where
  Top-20 alone achieves 0\%.
\item Clean accuracy degradation is $<$4\% for most modulations
  (worst: QAM16 at $-$8\%).
\item The defense is attack-agnostic: identical parameters work
  for both $L_2$ and $L_1$ attacks.
\end{itemize}

\subsection{Comparison with Alternative Defenses}

Table~\ref{tab:17def} compares representative defenses at SNR~18~dB
under EADL1. The spectral-gated defense achieves the highest average
while being the only method to recover AM-SSB.

\begin{table}[t]
\caption{Defense Comparison at SNR 18 dB (EADL1). Accuracy (\%)
  per modulation. Only Gated recovers all modulations.}
\label{tab:17def}
\centering
\scriptsize
\begin{tabular}{lccccc}
\toprule
Mod & Top-20 & LP-20 & Quant32 & MvAvg3 & \textbf{Gated} \\
\midrule
BPSK   & 77.0 & 74.0 &  44.5 & 88.5 & \textbf{77.0} \\
QPSK   & 84.5 & 89.0 &  93.0 & 86.5 & \textbf{84.5} \\
8PSK   & 54.5 & 70.0 &  89.5 & 17.5 & \textbf{54.0} \\
QAM16  & 59.5 & 62.5 &  51.5 & 56.5 & \textbf{59.5} \\
QAM64  & 89.5 & 85.5 &  61.5 & 84.5 & \textbf{89.5} \\
PAM4   & 95.5 & 95.5 &  93.0 & 95.0 & \textbf{95.5} \\
AM-DSB & 98.0 & 95.0 &  43.0 & 96.0 & \textbf{98.0} \\
GFSK   & 95.5 & 94.5 &  71.5 & 95.5 & \textbf{95.5} \\
CPFSK  & 72.5 & 85.0 &  77.5 & 68.0 & \textbf{72.5} \\
AM-SSB &  0.0 &  0.0 &  70.5 &  0.0 & \textbf{65.5} \\
WBFM   & 40.0 & 38.0 &  27.0 & 37.0 & \textbf{40.0} \\
\midrule
AVG    & 69.7 & 71.7 &  65.7 & 65.9 & \textbf{75.6} \\
\bottomrule
\end{tabular}
\end{table}

No individual defense dominates: Top-20 excels for QAM64/PAM4 but
fails on AM-SSB; Quant32 saves AM-SSB but destroys BPSK/AM-DSB;
MvAvg3 is best for BPSK but worst for 8PSK. The spectral-gated
defense matches Top-20 on narrowband modulations (Diff = 0\%) while
gaining 65.5\% on AM-SSB via quantization routing.

\subsection{Receiver Pipeline Analysis}

A natural question is whether standard receiver-front-end blocks
(AGC, matched filter) provide implicit defense. We evaluate these
at SNR~18~dB under EADL1 (Table~\ref{tab:rxpipe}).

\begin{table}[t]
\caption{Receiver Pipeline Blocks as Defense (EADL1, SNR 18 dB).
  Standard blocks degrade accuracy because the AWN model was trained
  on raw IQ signals, not preprocessed ones.}
\label{tab:rxpipe}
\centering
\small
\begin{tabular}{lc}
\toprule
Method & Avg Accuracy (\%) \\
\midrule
No defense (raw adversarial) & 0.0 \\
AGC only & 12.9 \\
DC removal + AGC & 10.2 \\
Matched filter (RRC) & 17.3 \\
AGC + matched filter & 12.4 \\
Resampling (2$\times$) & 64.2 \\
\midrule
\textbf{Spectral-Gated Defense} & \textbf{75.6} \\
\bottomrule
\end{tabular}
\end{table}

AGC and matched filtering \emph{reduce} accuracy to 10--17\%
because the AWN model was trained on raw IQ samples. Normalizing
amplitude (AGC) or filtering pulse shape (MF) alters the feature
distribution the model expects. Only resampling (which acts as
implicit lowpass filtering) provides meaningful recovery (64.2\%),
but still below the proposed defense.

%-------------------------------------------------------------------------------
\section{FEC on the Data Path}
\label{sec:fec}
%-------------------------------------------------------------------------------

The spectral-gated defense recovers \emph{classification accuracy},
enabling the receiver to select the correct demodulator. However,
defense filtering introduces residual signal distortion that can
cause symbol errors during demodulation. We show that forward error
correction (FEC) on the data path is complementary.

\subsection{FEC Configuration}

We use a rate-1/2 convolutional code with constraint length $K{=}7$
(generators [171, 133] octal), 64-state Viterbi decoder with soft
LLR inputs, and a 14-column block interleaver.

\subsection{FEC Results}

Table~\ref{tab:fec} compares CRC pass rates with and without FEC
after defense filtering at SNR~18~dB under CW attack, using oracle
demodulation.

\begin{table}[t]
\caption{CRC Pass Rate (\%) With/Without FEC After Defense (CW,
  SNR 18 dB, Oracle Demod). FEC dramatically improves data
  integrity after defense filtering.}
\label{tab:fec}
\centering
\small
\begin{tabular}{l|cc|cc|cc}
\toprule
 & \multicolumn{2}{c|}{No Def} & \multicolumn{2}{c|}{Top-20}
 & \multicolumn{2}{c}{LP-20} \\
Mod & noF & FEC & noF & FEC & noF & FEC \\
\midrule
QPSK  & 100 & 100 & 73.5 & \textbf{100} & 75.5 & \textbf{100} \\
8PSK  & 100 & 100 &  8.0 & \textbf{100} &  7.0 & \textbf{100} \\
QAM16 & 99.5 & 100 &  2.0 & \textbf{90} &  1.0 & \textbf{91} \\
QAM64 & 74.5 & \textbf{99} &  1.0 & 23 &  0.5 & 24 \\
PAM4  & 68.0 & \textbf{93.5} &  6.0 & \textbf{84.5} &  5.5 & \textbf{87} \\
\bottomrule
\end{tabular}
\end{table}

FEC transforms near-total CRC failure into high integrity:
8PSK goes from 8\% (Top-20, no FEC) to 100\% (Top-20 + FEC).
Even without any defense, FEC improves QAM64 from 74.5\% to 99\%.

\subsection{Recommended Architecture}

Based on our findings, we propose a split-path receiver architecture:

\begin{equation*}
\text{Signal} \rightarrow
\begin{cases}
\text{Spectral-Gated} \rightarrow \text{AMC} & \text{(control path)} \\
\text{Original signal} \rightarrow \text{Demod} \rightarrow \text{FEC} & \text{(data path)}
\end{cases}
\end{equation*}

The defense filter is applied only to the \emph{classification path}.
The data path uses the original (unfiltered) signal with FEC decoding,
avoiding defense-induced distortion on the demodulated data.

%-------------------------------------------------------------------------------
\section{Discussion}
\label{sec:discuss}
%-------------------------------------------------------------------------------

\subsection{Why Not Ensemble Voting?}

An alternative approach is to apply multiple defenses, classify each,
and vote on the result. While this achieves slightly higher accuracy
for some modulations, it requires $M$ classifier inferences per sample
(e.g., 3 for a 3-defense ensemble). At 35~ms per inference, this
violates real-time constraints for high-throughput receivers. The
spectral-gated defense achieves comparable accuracy with exactly
one inference.

\subsection{Limitations}

\textbf{AM-SSB recovery gap.} While the defense improves AM-SSB from
0\% to 20--65\%, this remains below the 93.5\% clean accuracy. Input
quantization is a coarse defense for wideband signals, and improved
wideband-aware methods are needed.

\textbf{Low-SNR regime.} At SNR~0~dB, Top-20 removes too many
signal components, degrading clean accuracy for BPSK (54.5\%), QPSK
(22.5\%), and 8PSK (2\%). A larger $K$ or SNR-adaptive selection
would mitigate this at the cost of reduced attack filtering.

\textbf{Model dependence.} The defense is designed for the raw-IQ
input distribution that the AWN model expects. Models trained with
different preprocessing (e.g., post-AGC inputs) would require
different defense parameterization.

%-------------------------------------------------------------------------------
\section{Related Work}
\label{sec:related}
%-------------------------------------------------------------------------------

\textbf{Adversarial attacks on AMC.}
Flowers et al.~\cite{flowers2019sigguard} evaluated FGSM and PGD
attacks on CNN-based AMC and proposed SigGuard, a detector-based
defense. Lin et al.~\cite{lin2020tactics} demonstrated universal
perturbation attacks on DNN-based modulation recognition.
Silvaco et al.~\cite{silvaco2023adversarial_rf} survey the
attack landscape. Our work extends this by characterizing the
attacks as control-plane and proposing input-transformation defenses
rather than detector-based approaches.

\textbf{Input transformation defenses.}
In the image domain, input transformations (JPEG compression,
spatial smoothing, bit-depth reduction) have been proposed as
pre-processing defenses. We adapt this concept to IQ signals,
finding that FFT Top-$K$ filtering is the RF-domain analog of
spatial frequency filtering, and input quantization corresponds
to bit-depth reduction.

\textbf{FEC and adversarial robustness.}
While FEC is standard in digital communications~\cite{proakis2008digital},
its interaction with adversarial attacks on AMC has not been
studied. We show that FEC addresses a fundamentally different
failure mode (residual symbol errors) than the defense
(classification recovery).

%-------------------------------------------------------------------------------
\section{Conclusion}
\label{sec:conclusion}
%-------------------------------------------------------------------------------

We demonstrate that adversarial attacks on deep-learning-based AMC
are control-plane attacks: they manipulate the classification
decision without corrupting the physical signal. Exploiting this
insight, we propose Spectral-Gated Defense, a real-time
input-transformation method that uses spectral flatness to route
signals to modulation-appropriate filtering. Combined with FEC on
the data path, the defense provides a practical, deployable
solution for robust AMC systems. All 11 modulations in RML2016.10a
show positive recovery under both CW and EADL1 attacks, with
sub-millisecond per-sample latency.

%-------------------------------------------------------------------------------
\section*{Availability}
%-------------------------------------------------------------------------------

The implementation of \textsc{Spectral-Gated Defense} is available
as \texttt{spectral\_gated\_defense()} in \texttt{util/defense.py}.
Experiments are reproducible using the scripts and checkpoints in
the accompanying repository.

%-------------------------------------------------------------------------------
\bibliographystyle{plain}
\bibliography{\jobname}

%%%%%%%%%%%%%%%%%%%%%%%%%%%%%%%%%%%%%%%%%%%%%%%%%%%%%%%%%%%%%%%%%%%%%%%%%%%%%%%%
\end{document}
%%%%%%%%%%%%%%%%%%%%%%%%%%%%%%%%%%%%%%%%%%%%%%%%%%%%%%%%%%%%%%%%%%%%%%%%%%%%%%%%
